\documentclass[aspectratio=169,11pt]{beamer}

\usetheme{Madrid}
\usecolortheme{default}
\usefonttheme{professionalfonts}
\setbeamertemplate{navigation symbols}{}
\setbeamertemplate{footline}[frame number]

\usepackage{amsmath, amssymb, booktabs, graphicx}

\title{Worker Voice, Collective Action, and Bargaining Institutions}
\subtitle{Labor Markets and Political/Cultural Institutions --- Week 4}
\author{Teaching deck draft}
\date{}

\begin{document}

\begin{frame}
  \titlepage
\end{frame}

\begin{frame}{Central question}
\begin{itemize}
    \item How do worker organizations and bargaining institutions form?
    \item How do they change wages, rents, governance, inequality, and workplace power?
    \item Which effects are direct effects on organized units, and which are spillovers to uncovered workers or politics?
\end{itemize}
\end{frame}

\begin{frame}{Bridge from Week 3}
\begin{itemize}
    \item Week 3: formal contracts, payroll systems, and legal protection often cover only part of work.
    \item Week 4: workers may respond to weak individual claim power through collective organization.
    \item Collective institutions change threat points, information flows, monitoring, and surplus division.
    \item The labor object is organization, coverage, wage-setting, voice, and spillovers.
\end{itemize}
\end{frame}

\begin{frame}{What counts as worker voice?}
\small
\begin{tabular}{p{0.28\linewidth} p{0.58\linewidth}}
\toprule
Object & Labor-market interpretation \\
\midrule
Union membership & workers formally join an organization \\
Bargaining coverage & workers are covered by a collective agreement \\
Enterprise bargaining & bargaining occurs at the firm or establishment \\
Sectoral bargaining & bargaining is coordinated above the firm \\
Works councils & shop-floor consultation and representation \\
Codetermination & board-level worker representation \\
Strike capacity & credible collective-action threat \\
\bottomrule
\end{tabular}
\end{frame}

\begin{frame}{Membership is not coverage}
\begin{itemize}
    \item Membership measures the workers who join an organization.
    \item Coverage measures the workers governed by a collective agreement.
    \item Coverage can exceed membership when sectoral agreements or extension rules apply.
    \item Enterprise systems make local certification and employer opposition central.
    \item Sectoral systems make coordination, extension, and wage compression central.
\end{itemize}
\end{frame}

\begin{frame}{How organizations form}
\small
\begin{tabular}{p{0.26\linewidth} p{0.30\linewidth} p{0.30\linewidth}}
\toprule
Margin & Worker side & Employer / legal side \\
\midrule
Demand & voice, wages, fairness, protection & expected cost of organization \\
Costs & recruitment, trust, dues, retaliation risk & opposition and delay \\
Certification & votes, mobilization, majority thresholds & recognition law and unit rules \\
Tightness & outside options and replacement risk & labor scarcity and vacancy costs \\
Design & union, works council, sectoral coverage & extension, thresholds, strike law \\
\bottomrule
\end{tabular}
\end{frame}

\begin{frame}{Organization is endogenous}
\begin{itemize}
    \item Observed union density is not latent worker demand.
    \item It is demand filtered through organizing costs, law, employer opposition, and labor-market conditions.
    \item Close elections identify marginal recognition victories, not all possible unionization.
    \item Tight labor markets can raise activity by improving outside options and lowering retaliation fear.
    \item Institutional design determines whether organization is workplace-by-workplace or coverage-based.
\end{itemize}
\end{frame}

\begin{frame}{Wages, rents, and coverage}
\begin{itemize}
    \item Collective bargaining can shift rents from firms to workers.
    \item Wage effects depend on firm rents, labor-demand elasticity, product-market competition, and strike capacity.
    \item Coverage can matter even when membership is low.
    \item Sectoral bargaining can compress wages across workers and firms.
    \item Enterprise bargaining is closest to close-election causal designs.
\end{itemize}
\end{frame}

\begin{frame}{A surplus map}
\[
\Delta S
= \Delta W_{workers}
+ \Delta \Pi_{firms}
+ \Delta G_{voice}
+ \Delta E_{spillovers}
\]
\begin{itemize}
    \item $\Delta W_{workers}$: wages, benefits, risk, and job quality for covered workers.
    \item $\Delta \Pi_{firms}$: profits, productivity, payroll, investment, and employment.
    \item $\Delta G_{voice}$: governance, information, retention, and representation.
    \item $\Delta E_{spillovers}$: uncovered workers, wage norms, inequality, and politics.
\end{itemize}
\end{frame}

\begin{frame}{Close-election evidence}
\begin{itemize}
    \item DiNardo--Lee: new unionization near certification-election thresholds.
    \item Frandsen: matched employer-employee evidence on worker and establishment margins.
    \item The design compares barely winning and barely losing campaigns.
    \item It is strongest for direct effects on organized units near the threshold.
    \item It is not a complete estimate of sectoral coverage, spillovers, or latent demand.
\end{itemize}
\end{frame}

\begin{frame}{Voice and governance}
\begin{itemize}
    \item The voice face changes information, grievance handling, scheduling, training, safety, and retention.
    \item Works councils provide consultation and workplace representation.
    \item Codetermination gives workers board-level representation.
    \item These institutions can affect productivity and trust even when wage bargaining is limited.
    \item The empirical object is governance, not just wage premia.
\end{itemize}
\end{frame}

\begin{frame}{German model as a teaching case}
\begin{itemize}
    \item Sectoral bargaining can set wage standards above the firm.
    \item Works councils can govern workplace information and consultation.
    \item Codetermination can connect labor voice to corporate decisions.
    \item The same country can combine multiple representation channels.
    \item This is why ``union effect'' needs institutional detail.
\end{itemize}
\end{frame}

\begin{frame}{Spillovers and inequality}
\begin{itemize}
    \item Threat effects: nonunion firms may raise wages to deter organization.
    \item Coverage spillovers: agreements can influence uncovered workers.
    \item Norm spillovers: wage-setting rules can create broader pay standards.
    \item Composition effects: firms may change hiring, outsourcing, or capital use.
    \item Inequality effects require direct and uncovered-worker margins.
\end{itemize}
\end{frame}

\begin{frame}{Direct effects are not enough}
\begin{itemize}
    \item A treatment-on-the-treated estimate can miss wage norms in the wider market.
    \item Coverage can compress wages even when membership is modest.
    \item Uncovered workers may gain through threat effects or lose through reallocation.
    \item Distributional decompositions ask how institutions reshape the whole wage structure.
    \item Fortin--Lemieux--Lloyd is the Week 4 spillover anchor.
\end{itemize}
\end{frame}

\begin{frame}{Political spillovers and feedback}
\small
\begin{tabular}{p{0.30\linewidth} p{0.55\linewidth}}
\toprule
Direction & Mechanism \\
\midrule
Unions to politics & turnout, political information, participation, campaign contact, policy advocacy \\
Politics to unions & right-to-work laws, bargaining rights, strike law, public-sector rules, enforcement \\
Feedback loop & political rules reshape worker organization; organizations reshape political coalitions \\
\bottomrule
\end{tabular}
\end{frame}

\begin{frame}{Keep politics bounded}
\begin{itemize}
    \item Political outcomes matter because unions are worker organizations.
    \item The labor-market object remains organization, coverage, bargaining power, and voice.
    \item Voting and participation are downstream outcomes or feedback channels.
    \item Right-to-work laws are institutional shocks to organizing capacity and union security.
    \item This is a labor institutions lecture, not a general parties lecture.
\end{itemize}
\end{frame}

\begin{frame}{Comparative regimes}
\small
\begin{tabular}{p{0.26\linewidth} p{0.28\linewidth} p{0.32\linewidth}}
\toprule
Regime & Central object & Main labor margin \\
\midrule
Enterprise bargaining & local recognition & firm-level wages and employment \\
Sectoral bargaining & coverage and extension & wage coordination and compression \\
Dual-channel systems & bargaining plus councils & wages, governance, information \\
Codetermination & board representation & governance and strategy \\
Weak coverage & high costs or legal barriers & latent demand and failed formation \\
\bottomrule
\end{tabular}
\end{frame}

\begin{frame}{Empirical designs and what they identify}
\small
\begin{tabular}{p{0.34\linewidth} p{0.25\linewidth} p{0.27\linewidth}}
\toprule
Design & Identifies best & Observed margins \\
\midrule
Close-election RD & recognition at threshold & wages, employment, payroll \\
Matched event study & worker and firm incidence & careers, separations, turnover \\
Bargaining-right reform & legal capacity to organize & coverage, wages, politics \\
Representation threshold & voice institutions & productivity, retention, governance \\
Distributional decomposition & spillovers and inequality & wage structure, uncovered workers \\
Survey / field evidence & organizing demand & preferences, fear, tightness \\
\bottomrule
\end{tabular}
\end{frame}

\begin{frame}{Week 4 lab}
\begin{itemize}
    \item Reproduce: synthetic close-election design inspired by DiNardo--Lee.
    \item Diagnose: worker and establishment margins inspired by Frandsen.
    \item Transfer: classify spillover, governance, and political-feedback designs.
    \item Core output: what object forms, what margin moves, and what the design identifies.
\end{itemize}
\end{frame}

\begin{frame}{Bridge to Week 5}
\begin{itemize}
    \item Week 4: formal collective organization changes bargaining, governance, and spillovers.
    \item Week 5: culture and norms shape labor allocation before formal organization may exist.
    \item The next question: how do beliefs about work, family, status, mobility, and acceptable jobs shape labor supply and search?
    \item Keep the same discipline: labor outcome, institutional object, mechanism, evidence.
\end{itemize}
\end{frame}

\begin{frame}{Takeaways}
\begin{itemize}
    \item Collective institutions are about organization, coverage, bargaining, voice, and power.
    \item Organization formation is endogenous to worker demand, law, costs, opposition, and tightness.
    \item Wage effects, governance effects, spillovers, and political feedback are distinct margins.
    \item Strong empirical work names the institutional object and the identified margin.
\end{itemize}
\end{frame}

\end{document}
